\chapter{Coding}

\section{Backend Entry Point (main.py)}
\begin{verbatim}
from fastapi import FastAPI
from fastapi.middleware.cors import CORSMiddleware
from routers.auth      import router as auth_router
from routers.ids       import router as ids_router
from routers.phishing  import router as phishing_router
from routers.risk      import router as risk_router
from routers.dashboard import router as dashboard_router
from routers.notifications import router as notifications_router
from middleware.request_logger import RequestLoggerMiddleware

app = FastAPI(
    title="SmartSec",
    description="AI-Based Cyber Defense Platform",
    version="1.0.0",
)
app.add_middleware(CORSMiddleware, allow_origins=ALLOWED_ORIGINS,
                   allow_credentials=True, allow_methods=["*"],
                   allow_headers=["*"])
app.add_middleware(RequestLoggerMiddleware)
app.include_router(auth_router)
app.include_router(ids_router)
app.include_router(phishing_router)
app.include_router(risk_router)
app.include_router(dashboard_router)
app.include_router(notifications_router)
\end{verbatim}
\textit{The application is launched using Uvicorn ASGI server with hot reload during development. Six modular routers cover all security modules. A custom RequestLoggerMiddleware intercepts every incoming request to feed the IDS anomaly detection pipeline.}

\section{Authentication Module (auth.py)}
\begin{verbatim}
@router.post("/register", response_model=Token,
             status_code=status.HTTP_201_CREATED)
async def register(payload: UserRegister,
                   db: Client = Depends(get_supabase)):
    # 1. Check for duplicate email and username
    existing = db.table("users").select("id") \
                 .eq("email", payload.email).execute()
    if existing.data:
        raise HTTPException(409, "Email already exists.")
    # 2. Hash password with bcrypt (12 rounds)
    hashed = hash_password(payload.password)
    # 3. Insert new user into Supabase
    result = db.table("users").insert({
        "username": payload.username,
        "email": payload.email,
        "hashed_password": hashed,
        "risk_score": 0.0, "risk_level": "Low",
    }).execute()
    # 4. Issue JWT token
    token = create_access_token({"sub": result.data[0]["email"]})
    return Token(access_token=token, token_type="bearer", ...)
\end{verbatim}
\textit{The auth router implements registration, login, OAuth callback, profile update, settings persistence, avatar upload to Supabase Storage, password change, and login history endpoints. Login risk assessment evaluates off-hours access, new IP detection, and failed-attempt count to assign risk flags.}

\section{Phishing Detection Service}
\begin{verbatim}
@router.post("/scan")
async def scan_url(body: ScanRequest,
                   user=Depends(get_current_user),
                   db=Depends(get_supabase)):
    result = analyse_url(body.url)  # 19 heuristics + 3 APIs

    # Persist to Supabase url_scans table
    db.table("url_scans").insert({
        "user_id":  user["id"], "url": body.url,
        "verdict":  result.verdict,
        "confidence": result.confidence,
        "risk_score": result.risk_score,
        "features": result.features,
    }).execute()

    # Trigger risk engine side-effect
    if result.verdict == "Phishing":
        record_event(db, user, "phishing_detected",
          description=f"Phishing URL: {body.url[:100]}")
    elif result.verdict == "Suspicious":
        record_event(db, user, "phishing_suspicious",
          description=f"Suspicious URL: {body.url[:100]}")
    return {"verdict": result.verdict,
            "risk_score": result.risk_score, ...}
\end{verbatim}
\textit{The phishing service evaluates 19 URL signals (length, entropy, IP-based URL, suspicious TLD, excessive subdomains, URL shortener, at-symbol, double-slash, hex encoding, etc.) and queries VirusTotal, Google Safe Browsing, and AbuseIPDB. A weighted aggregate of heuristic and API signals produces a 0--100 risk score with a Safe/Suspicious/Phishing verdict.}

\section{Risk Scoring Engine (risk\_service.py)}
\begin{verbatim}
EVENT_WEIGHTS = {
    "login_failure":       (5.0,  "Low"),
    "suspicious_login":    (12.0, "Medium"),
    "phishing_suspicious": (10.0, "Medium"),
    "phishing_detected":   (22.0, "High"),
    "ids_anomaly_low":     (5.0,  "Low"),
    "ids_anomaly_high":    (20.0, "High"),
    "brute_force":         (28.0, "Critical"),
    "sql_injection":       (22.0, "High"),
}
DECAY_PER_DAY = 0.08  # 8% exponential daily decay

def record_event(db, user, event_type, description="",
                 metadata=None, ip=None):
    delta, severity = EVENT_WEIGHTS.get(event_type,(5.0,"Low"))
    current = float(user.get("risk_score", 0) or 0)
    # Apply time decay before adding new delta
    current = _apply_decay(current, user.get("updated_at"))
    new_score = min(round(current + delta, 2), 100.0)
    new_level = score_to_level(new_score)
    # Persist event + update user + create notification
    db.table("risk_events").insert({...}).execute()
    db.table("users").update({
        "risk_score": new_score, "risk_level": new_level
    }).eq("id", user["id"]).execute()
    # Email alert for High/Critical
    if severity in ("High", "Critical") and user.get("email"):
        send_alert_email(to=user["email"], ...)
    return {"new_score": new_score, "new_level": new_level}
\end{verbatim}
\textit{The risk engine maintains a 0--100 per-user score. Each security event applies a weighted delta based on event severity. An 8\% per-day exponential decay is applied to the existing score before adding the new delta, ensuring old incidents gradually fade. Risk levels are: Low (0--24), Medium (25--49), High (50--74), Critical (75--100).}

\section{IDS Service (IsolationForest)}
\begin{verbatim}
from sklearn.ensemble import IsolationForest
import numpy as np

FEATURE_NAMES = ["packet_size","protocol_type","request_rate",
                 "port","duration","error_rate","bandwidth"]

class IDSModel:
    def __init__(self):
        self.model = IsolationForest(
            n_estimators=100, contamination=0.1,
            random_state=42)
        self._train_on_synthetic()

    def _train_on_synthetic(self):
        np.random.seed(42)
        normal = np.random.randn(1000, len(FEATURE_NAMES))
        anomalies = np.random.randn(100, len(FEATURE_NAMES))
        anomalies[:, 2] *= 5   # spike request_rate
        anomalies[:, 5] *= 4   # spike error_rate
        X = np.vstack([normal, anomalies])
        self.model.fit(X)

    def predict(self, features: dict) -> dict:
        x = np.array([[features[f] for f in FEATURE_NAMES]])
        pred = self.model.predict(x)[0]
        score = -self.model.score_samples(x)[0]
        return {"is_anomaly": pred == -1,
                "anomaly_score": round(float(score), 4)}
\end{verbatim}
\textit{The IDS module uses scikit-learn's IsolationForest unsupervised algorithm trained on synthetic traffic data at startup. The model isolates anomalous data points using random partitioning trees. A contamination parameter of 0.1 (10\%) is set to reflect the expected proportion of anomalous traffic. The anomaly score (normalised decision function output) provides a continuous confidence measure.}

\section{Notification and Dashboard Modules}
\begin{verbatim}
# GET /notifications
@router.get("")
async def list_notifications(user=Depends(get_current_user),
                             db=Depends(get_supabase)):
    res = db.table("notifications").select("*") \
             .eq("user_id", user["id"]) \
             .order("created_at", desc=True) \
             .limit(50).execute()
    unread = sum(1 for n in res.data if not n["is_read"])
    return {"notifications": res.data, "unread_count": unread}

# GET /dashboard
@router.get("")
async def dashboard(user=Depends(get_current_user),
                    db=Depends(get_supabase)):
    # Aggregate data from all 4 security modules
    phishing_stats  = get_phishing_summary(db, user["id"])
    ids_stats       = get_ids_summary(db, user["id"])
    risk_data       = get_risk_summary(db, user)
    activity_data   = get_login_activity(db, user["id"])
    return {"phishing": phishing_stats, "ids": ids_stats,
            "risk": risk_data, "activity": activity_data}
\end{verbatim}
\textit{The notification centre stores per-user alert records in Supabase and returns an unread count on every poll. The React frontend polls the /notifications endpoint every 30 seconds using a context-level interval. The dashboard endpoint aggregates live data from all four security modules into a single response to minimise frontend API calls.}
